\documentclass{article}

% ---- Packages ----
\usepackage[preprint]{neurips_2024}   % swap to [final] for camera-ready
\usepackage[utf8]{inputenc}
\usepackage[T1]{fontenc}
\usepackage{hyperref}
\usepackage{url}
\usepackage{booktabs}
\usepackage{amsmath}
\usepackage{amssymb}
\usepackage{graphicx}
\usepackage{microtype}
\usepackage{xcolor}
\usepackage{enumitem}

% ---- Convenience macros ----
\newcommand{\REQ}{\texttt{[REQ]}}
\newcommand{\TEST}{\texttt{[TEST]}}
\newcommand{\ORIG}{\texttt{[ORIG]}}
\newcommand{\HUNK}{\texttt{[HUNK]}}
\newcommand{\tier}[1]{\textbf{T#1}}
\newcommand{\todo}[1]{\textcolor{red}{[TODO: #1]}}

% -----------------------------------------------------------------------
\title{The Geometry of Edit Necessity:\\
Learning What a Requirement Change Entails}

\author{
  Yundong Shen \quad Miaowen Dong \\
  \textit{Chalmers University of Technology} \\
  \texttt{\{yundong, miaowen\}@chalmers.se}
}

\begin{document}
\maketitle

% -----------------------------------------------------------------------
\begin{abstract}
Pre-trained code models capture semantic similarity between natural language
and code, yet they cannot reliably judge whether a specific code change is
\emph{necessary} given a requirement---even when it is semantically relevant.
We ask whether edit necessity is a geometrically learnable property of a joint
embedding space.
We propose a self-supervised contrastive framework that jointly embeds four
entities---requirement (\REQ), tests (\TEST), source code units (\ORIG), and
diff hunks (\HUNK)---trained on SWE-bench Full (2{,}291 instances, 12
repositories) without any human-provided necessity labels.
We report three findings.
First, necessity tiers form statistically separable regions in the learned
space (silhouette score rising from $0.010$ to $0.139$ after training).
Second, edit necessity decomposes into two orthogonal components: a
\emph{semantic} component already partially captured by pre-trained models,
and a \emph{structural} component---whether the modified code unit is
implicated by the requirement---that requires the \ORIG{} entity and is
invisible to text-only representations.
Third, this geometric structure does not emerge from standard InfoNCE over
random negatives; it requires hard negatives drawn from LLM-generated
scope-creep hunks that are semantically plausible but structurally
unnecessary.
On a held-out test set, our framework recovers $92.4\%$ of maintainer-endorsed
but test-invisible changes (Tier-2 recall), compared to $80.6\%$ for an
untrained UniXCoder baseline, with the gap attributable almost entirely to the
structural component.
\end{abstract}

% -----------------------------------------------------------------------
\section{Introduction}
\label{sec:intro}

Large pre-trained code models have demonstrated strong performance across a
wide range of tasks that require semantic understanding: code search, clone
detection, summarization, and defect prediction
\citep{feng2020codebert,guo2022unixcoder}.
These models learn rich representations of what code \emph{means}, and they
transfer well to tasks that reduce to semantic similarity.
Yet there is a class of judgment that these representations do not support:
given a requirement change, determining which code modifications are
\emph{necessary}---and which, while plausible, go beyond what the requirement
entails.

This distinction matters.
Two diff hunks can be equally related to a bug report in the space of text;
one may be indispensable to the fix while the other introduces unrelated
behaviour.
The difference does not lie in what the code \emph{says}, but in what the
requirement \emph{entails} given the structure of the existing codebase and
the boundaries of the test suite.
Neither property is encoded in pre-trained code representations.

\paragraph{An implicit signal in code evolution.}
Human developers make necessity judgments implicitly with every commit.
When a maintainer merges a patch for a bug report, they determine through code
review and domain knowledge which changes belong to the fix and which do not.
This judgment is shaped jointly by the requirement's semantics, the test
suite's coverage, and the structural relationships within the codebase.
Large repositories of resolved issues---such as SWE-bench
\citep{jimenez2024swebench}---encode thousands of such decisions.
Yet no existing approach learns the \emph{geometry} of necessity from this
signal.

\paragraph{Research question.}
We ask: \emph{is edit necessity geometrically learnable in a joint embedding
space?}
More precisely, given an embedding function trained on (requirement, test,
source, hunk) quadruples, do hunks at different necessity levels occupy
statistically distinguishable regions, and does the model's geometric
structure support recovering necessary changes that are invisible to tests?

\paragraph{Approach.}
We introduce \textbf{Edit Entailment Learning} (EEL), a self-supervised
contrastive framework that embeds four entities into a shared $\ell_2$-normalised
256-dimensional space.
Type tokens (\REQ, \TEST, \ORIG, \HUNK) condition a shared UniXCoder
\citep{guo2022unixcoder} backbone, and a MultiPair InfoNCE loss pulls four
types of positive pairs together while pushing apart in-batch negatives.
We define three necessity tiers derived entirely from SWE-bench's structure,
requiring no human annotation:
\begin{itemize}[itemsep=2pt,topsep=2pt]
  \item \tier{1}: gold hunks covered by fail-to-pass tests (highest necessity)
  \item \tier{2}: gold hunks \emph{not} covered by any test (maintainer-endorsed
    but test-invisible)
  \item \tier{3}: LLM-generated hunks outside the gold patch (scope creep)
\end{itemize}
Tier-2 is the critical case: these changes are necessary for correctness, yet
every test-based method---static analysis, coverage-guided review,
execution-based repair---treats them as untouched and therefore unverified.

\paragraph{Findings.}
We report three findings that together characterise the geometry of edit
necessity.

\textbf{Finding 1: Necessity is geometrically separable, but only after
training with targeted hard negatives.}
An untrained UniXCoder baseline yields a cosine-space silhouette score of
$0.010$ across necessity tiers---effectively random.
Adding basic InfoNCE over random in-batch negatives (M1) changes this
negligibly ($0.011$).
Only when LLM-generated scope-creep hunks serve as hard negatives (M2) does
the score rise to $0.139$, with Mann-Whitney $U$ confirming significant
separation at every tier boundary ($p < 10^{-25}$ for all gold-vs-drift
comparisons).
The finding implies that semantic hard negatives of a specific kind---
plausible but structurally unnecessary---are necessary to induce geometric
structure, a distinction that random sampling cannot provide.

\textbf{Finding 2: Edit necessity decomposes into two orthogonal components.}
The untrained model already achieves $80.6\%$ Tier-2 recall, which at first
suggests the task is easy.
We show this reflects a \emph{semantic} component of necessity---hunks whose
code is textually related to the requirement---that UniXCoder's pre-training
already captures.
The remaining gap ($80.6\% \to 92.4\%$) is a \emph{structural} component:
hunks that modify code units structurally implicated by the requirement, even
when the connection is not explicit in the text.
We operationalise this split via a Tier-2a/2b partition: T2a hunks modify
existing code units (high structural necessity) and T2b hunks add new code
(primarily semantic necessity).
The \ORIG{} entity, which encodes the source unit being modified, is the sole
carrier of structural necessity in the model.

\textbf{Finding 3: The \ORIG{} entity is necessary and sufficient for the
structural component.}
An ablation that removes \ORIG{} from training (while retaining \REQ{} and
\TEST) drops Tier-2 recall back toward the baseline.
In contrast, removing \TEST{} leaves Tier-2 recall largely intact while
degrading Tier-1 ranking.
This dissociation directly supports the two-component decomposition: \TEST{}
encodes the verified boundary (relevant to T1), \ORIG{} encodes structural
implication (relevant to T2a), and neither is redundant with the other.
\todo{Fill in ablation numbers once jobs complete.}

\paragraph{Contributions.}
\begin{enumerate}[itemsep=2pt,topsep=2pt]
  \item We pose edit necessity as a geometric learning problem and provide the
    first empirical evidence that it is learnable from code evolution data
    without human annotation.
  \item We identify a two-component decomposition of necessity (semantic vs.\
    structural) and show that only the structural component requires learning;
    the semantic component is inherited from pre-training.
  \item We show that targeted hard negatives (structurally unnecessary but
    semantically plausible hunks) are a prerequisite for learning geometric
    necessity structure---a finding with implications for contrastive
    learning in structured domains.
\end{enumerate}

% -----------------------------------------------------------------------
\section{Background \& Problem Setup}
\label{sec:setup}
\todo{Write this section.}

% -----------------------------------------------------------------------
\section{Method}
\label{sec:method}
\todo{Write this section.}

% -----------------------------------------------------------------------
\section{Does the Geometry Exist?}
\label{sec:geometry}
\todo{Write this section.}

% -----------------------------------------------------------------------
\section{Two Components of Edit Necessity}
\label{sec:two-components}
\todo{Write this section.}

% -----------------------------------------------------------------------
\section{Retrieval Performance}
\label{sec:retrieval}
\todo{Write this section.}

% -----------------------------------------------------------------------
\section{Related Work}
\label{sec:related}
\todo{Write this section.}

% -----------------------------------------------------------------------
\section{Discussion \& Conclusion}
\label{sec:conclusion}
\todo{Write this section.}

% -----------------------------------------------------------------------
\bibliography{refs}
\bibliographystyle{plainnat}

\end{document}
